Hierarchical planning is a powerful approach for decision-making in environments with continuous states, continuous actions, and long horizons.
A crucial bottleneck in scaling hierarchical planning is the reliance on human engineers to manually program domain-specific abstractions.
For example, in bilevel sample-based task and motion planning~\cite{srivastava2014combined,garrett2021integrated}, an engineer must design (1) symbolic predicates; (2) symbolic operators; and (3) samplers that propose different refinements of the symbolic operators into continuous actions.
However, recent work has shown that when predicates are \emph{given}, operators and samplers can be learned from a modest number (50--200) of demonstrations~\cite{silver2021learning,chitnis2021learning}.
Our objective in this work is to \emph{learn} predicates that can then be used to learn operators and samplers.



\begin{figure*}[t]
  \centering
    \noindent
    \includegraphics[width=\textwidth]{figures/teaser.png}
    \caption{\small{\textbf{Overview of our framework.} Given a small set of goal predicates (first panel, top), we use demonstration data to learn new predicates (first panel, bottom). In this Blocks example, the learned predicates \texttt{P1} -- \texttt{P4} intuitively represent \texttt{Holding}, \texttt{NotHolding}, \texttt{HandEmpty}, and \texttt{NothingAbove} respectively. Collectively, the predicates define a state abstraction that maps continuous states $x$ in the environment to abstract states $s$. Object types are omitted for clarity. After predicate invention, we learn abstractions of the continuous action space and transition model via planning operators (second panel). For each operator, we learn a sampler (third panel), a neural network that maps continuous object features in a given state to continuous action parameters for controllers which can be executed in the environment. In this example, the sampler proposes different placements on the table for the held block.
    With these learned representations, we perform bilevel planning (fourth panel), with search in the abstract spaces guiding planning in the continuous spaces.}}
  \label{fig:teaser}
\end{figure*}

Predicates in bilevel planning represent a discrete state abstraction of the underlying continuous state space~\cite{li2006towards,abel2017near}.
For example, \texttt{On(block1, block2)} is an abstraction that discards the exact continuous poses of \texttt{block1} and \texttt{block2}.
State abstraction alone is useful for decision-making, but predicates go further: together with operators, predicates enable the use of highly-optimized domain-independent AI planners~\cite{helmert2006fast}.

We consider a problem setting where a small set of \emph{goal predicates} are available and sufficient for describing task goals, but practically insufficient for bilevel planning.
For example, in a block stacking domain~(\figref{fig:teaser}), we start with \texttt{On} and \texttt{OnTable}, but have no predicates for describing whether a block is currently held or graspable.
Our aim is to invent new predicates to enrich the state abstraction beyond what can be expressed with the goal predicates alone, leading to stronger reasoning at the abstract level.

What objective should we optimize to learn predicates for bilevel planning?
First, consider our real objective: we want a predicate set such that bilevel planning is fast and successful, in expectation over a task distribution, when we use those predicates to learn operators and samplers for planning.
Unfortunately, this real objective is far too expensive to use directly, since even a single evaluation requires neural network sampler training and bilevel planning.


In this work, we propose a novel surrogate objective that is deeply connected to our real bilevel-planning objective, but is tractable for predicate learning.
Our main insight is that demonstrations can be used to analytically approximate bilevel planning time.
To leverage this objective for predicate learning, we take inspiration from the program synthesis literature~\cite{menon2013machine,ellis2020dreamcoder}, and learn predicates via a hill-climbing search through a grammar, with the search guided by the objective.
After predicate learning, we use the predicates to learn operators and samplers.
All three components can then be used for efficient bilevel planning on new tasks.

In experiments across four robotic planning environments, we find predicates, operators, and samplers learned from 50--200 demonstrations enable efficient bilevel planning on held-out tasks that involve different numbers of objects, longer horizons, and larger goal expressions than seen in the demonstrations.
Furthermore, predicates learned with our proposed surrogate objective substantially outperform those learned with objectives inspired by previous work, which are based on prediction error~\cite{pasula2007learning,jetchev2013learning}, bisimulation~\cite{konidaris2018skills,curtis2021discovering}, and inverse planning~\cite{baker2009action,ramirez2010probabilistic,zhi2020online}.
We compare against several other baselines and ablations of our system to further validate our results.
